\documentclass[10pt]{article}
\usepackage[margin=20mm]{geometry}
\usepackage{amsmath,amssymb,mathtools,fontspec,unicode-math,xurl,hyperref}
\IfFontExistsTF{DejaVu Serif}{\setmainfont{DejaVu Serif}}{\setmainfont{Latin Modern Roman}}
\setmathfont{latinmodern-math.otf}
\setlength{\parindent}{0pt}\setlength{\parskip}{6pt}\sloppy
\title{Scientific initial review of
Syzygies of adjoint linear series on projective varieties}
\author{Assigned scientific review unit: M25-0148}
\date{Initial review; adjudication pending}
\begin{document}
\maketitle
\section*{Identity, scope, and assessment}
The paper is Syzygies of adjoint linear series on projective varieties, by Purnaprajna Bangere and Justin Lacini, associated with Duke Mathematical Journal 174 (2025), 473–499, DOI 10.1215/00127094-2024-0035 (https://doi.org/10.1215/00127094-2024-0035). The source actually reviewed is the 20-page author manuscript arXiv:2302.02517v1, dated 6 February 2023 [paper]. References to pages below mean its one-based PDF pages, agreeing with its printed pagination. The author publication list confirms the journal identity; the publisher article and public download endpoints did not provide the journal text. No equivalence with the journal proof is asserted.\par
Assessment: Major, a demonstrated failure of a core proof step. The claimed invariant isomorphism on p. 15 is false for the canonical Galois action in an example satisfying the main hypotheses. Branchwise transport of a canonical-form section need not glue at ramification. The manuscript uses this construction first for Lemma 5.10 and then for the higher-diagonal lifting in Theorems 5.11 and 5.12. This interrupts the displayed route from enlarged diagonals to the actual diagonals.\par
The example does not disprove Lemma 5.10 or either main theorem. In fact, unrestricted lifting is explicitly surjective in the example. The algebra of tree-shaped diagonal ideals and the enlarged-diagonal vanishing estimates have substantial positive checks, recorded below. The assessment concerns the specified author-version proof, not an assertion that its final results are false or irreparable.\par
The entire manuscript text was read. They are AI assistants within this initial review; no independent acceptance, human review, or independently verified different-model identity is claimed.\par
\section*{Operative definitions and dependency chain}
Throughout, \(X\) is smooth and projective over \(\mathbb C\), \(A\) is ample and globally generated, and \(L_m=\mathcal O_X(K_X+mA)\). Theorem 1.2 states \(N_p\) for \(m\ge n+1+p\); Theorem 1.4 states Koszulness for \(m\ge n+2\). Section 2 explicitly permits \(N_p\) for globally generated line bundles without requiring an embedding. Consequently \(X=\mathbb P^n,A=\mathcal O(1),m=n+1\), where \(L_m=\mathcal O\), is not a counterexample to the stated \(N_0\) convention.\par
A basepoint-free subspace of dimension \(n+1\) gives a finite flat map \(\alpha:X\to\mathbb P^n\). The enlarged diagonal is \(\Gamma_X=X\times_{\mathbb P^n}X\); it contains \(\Delta_X\), and its extra components meet the ordinary diagonal over the branch locus. The star arrangement \(\Delta^1\) uses edges \((1,j)\); the chain arrangement \(\Delta^2\) uses edges \((j-1,j)\). These are different ideals and serve different criteria.\par
The proof establishes vanishing for the enlarged arrangements in Theorem 5.8, with thresholds \(m\ge n+r-1\) for stars and \(m\ge n+2\) for chains. It then needs surjectivity of restrictions from enlarged to ordinary arrangements. The invariant construction on p. 15 and its reuse on p. 17 are the bridge to this required surjectivity; they are not consequences of Theorem 5.8 alone.\par
\section*{Problem and decisive evidence}
Problem 1 (Major: canonical Galois transport does not preserve branch gluing).  In Section 5.4, p. 15, let \(Y\) be the Galois closure and let \(G\) act on the second factor. For \[F=(\omega_Y\otimes\beta^*\mathcal O_X(mA))
       \boxtimes(\omega_Y\otimes\beta^*\mathcal O_X(mA)),\] the manuscript asserts that the injection into restrictions to the graph components induces an isomorphism \[(F|_{\Gamma_{Y_0}})^G
 \ \cong\
 \left(\bigoplus_{g\in G} F|_{g^*\Delta_{Y_0}}\right)^G.\] It deduces surjectivity onto the diagonal and subsequently extends sections by the Galois action in the p. 17 induction. The displayed invariant isomorphism fails at ramification.\par
\subsection*{An admissible global example}
Take \[X=Y=\mathbb P^1,\qquad
 \alpha:[z_0:z_1]\longmapsto[z_0^2:z_1^2],\qquad
 A=\alpha^*\mathcal O_{\mathbb P^1}(1)=\mathcal O_{\mathbb P^1}(2),\qquad m=2.\] Here \(n=1\), the space spanned by \(z_0^2,z_1^2\) is basepoint free, \(\alpha\) is finite flat and Galois, \(\beta\) is the identity, and \[L=\omega_{\mathbb P^1}\otimes\alpha^*\mathcal O_{\mathbb P^1}(2)=\mathcal O_{\mathbb P^1}(2).\] Thus all main assumptions and \(m\ge n+1\) hold. The involution is \(g(z)=-z\). The smooth locus is all of \(Y\); on a curve the required removal of a codimension-at-least-two set cannot discard the ramification points \(0,\infty\).\par
\subsection*{Local computation at a branch point}
On affine charts write the two coordinates as \(u,v\). The fiber product has ring \[R=\mathbb C[u,v]/(u^2-v^2).\] Its two graph components are \(v=u\) and \(v=-u\), meeting at \((0,0)\). The restriction injection is \[R\hookrightarrow \mathbb C[u]\oplus\mathbb C[u],\qquad
 h\longmapsto(h(u,u),h(u,-u)).\] Its image consists exactly of pairs with equal value at zero.\par
Let \(e\) be the local frame \(du\boxtimes dv\), including the frames of the pulled-back base line bundles. Those latter frames are invariant. The canonical second-factor action sends \(v\) to \(-v\) and \(e\) to \(-e\). Writing \(h=a(u)+v b(u)\), invariance of \(he\) requires \(a=0\). Its component values are therefore \[(u b(u),-u b(u))e.\] On the direct sum of graph restrictions, however, all pairs \((s(u),-s(u))e\) are invariant. The invariant pair \((1,-1)e\) is not in the restriction image, because \(1\ne-1\) at the node. Equivalently, the induced map on invariant modules is \[u\,\mathbb C[u]\ \hookrightarrow\ \mathbb C[u],\] with nonzero cokernel. Taking invariants is exact in characteristic zero; it does not force this cokernel’s invariant part to vanish.\par
This also identifies the obstruction to the proposed canonical identification of the two pulled-back canonical bundles by the ramification formula. The identification coming from the base differential has component values \(+1\) and \(-1\); it does not extend as a regular unit through the node. This statement does not exclude other, noncanonical isomorphisms of the abstract line bundles.\par
\subsection*{Global section count}
The discrepancy is visible in global sections as well. In the affine frame, \(H^0(L)\) consists of \(p(z)\,dz\) times the base frame, with \(\deg p\le2\), and \[g\cdot p(z)\,dz=-p(-z)\,dz.\] Its invariant subspace is one-dimensional, spanned by \(z\,dz\). Hence \(H^0(\mathbb P^1\times\mathbb P^1,F)^G\) has dimension \(3\).\par
The divisor \(\Gamma\) has bidegree \((2,2)\), while \(F=\mathcal O(2,2)\). The exact sequence \[0\longrightarrow\mathcal O(0,0)\longrightarrow\mathcal O(2,2)
       \longrightarrow F|_\Gamma\longrightarrow0\] has \(H^1(\mathcal O(0,0))=0\). Its one-dimensional kernel is anti-invariant, represented locally by \((u^2-v^2)\,du\,dv\). It follows that \[\dim H^0(\Gamma,F)^G=3,\qquad
 \dim\left(\bigoplus_{g\in G}H^0(g^*\Delta,F)\right)^G
       =h^0(\mathbb P^1,\mathcal O(4))=5.\] The first invariant space maps to the diagonal sections vanishing at both ramification points. In a degree-four polynomial basis these are the two independent endpoint conditions on the constant and leading coefficients. Exact polynomial and rational linear-algebra calculations in the verification record check the character dimensions and rank-two compatibility condition; the displayed module argument gives the general local reason.\par
\subsection*{Impact, counterevidence, and repair boundary}
The affected step supplies the Galois case of Lemma 5.10 (pp. 15–16). The general case then uses it in the trace diagram, and the p. 17 proof invokes the same Galois extension to patch higher unions. Thus the invalid isomorphism concerns the central derivation of both \(N_p\) and Koszulness. This explicit dependence, together with an admissible counterexample to the asserted map, supports the Major assessment.\par
There is decisive counterevidence against calling the final restriction claim false. For this example, \[H^0(\mathcal O(2))\otimes H^0(\mathcal O(2))\longrightarrow H^0(\mathcal O(4))\] is surjective: every monomial of degree at most four is a product of two monomials of degree at most two. A global section of \(F\) on the product restricts first to \(\Gamma\) and then to \(\Delta\). Therefore the unrestricted map \(H^0(\Gamma,F)\to H^0(\Delta,F)\) is surjective despite the invariant construction’s failure. The same polynomial multiplication verifies the required powers in this test case.\par
A repair must construct regular lifts satisfying the compatibility conditions on intersecting graph components, and must show that these lifts have the prescribed restrictions on the mixed unions used in the induction. Simply deleting the branch locus is insufficient because it is a divisor. Merely choosing another linearization in this cyclic example does not prove a uniform replacement compatible with the canonical trace construction. No complete general repair is certified here; no counterexample to Theorems 1.2 or 1.4 is asserted.\par
\section*{Positive checks and limits of coverage}
\paragraph*{Products, diagonal ideals, and trace.}\mbox{}\par
The separate-factor Tor vanishing in Lemmas 3.1–3.3 follows by flat extension over the ground field. The displayed identification \(A/\mathfrak m=\mathbb C\) should be restricted to the intended closed-point setting; the vanishing itself still follows with an arbitrary residue field, which is flat over \(\mathbb C\).\par
The star/chain statements in Corollary 3.5 admit a direct algebraic justification beyond the phrase change of coordinates. For a new leaf factor write \(R=B\otimes_\mathbb CA_{\rm new}\). Quotienting by its attaching diagonal gives \(B\), free over the old factor \(B\). A free resolution of an old ideal, extended to \(R\), stays exact after this quotient. The associated Tor vanishing gives product/intersection and tensor-product identifications inductively. Thus singular-coordinate concerns do not supply a counterexample to the forest-diagonal assertions. Likewise the normalized trace in characteristic zero splits \(\mathcal O_X\) from \(f_*\mathcal O_Y\); dualizing this split supports Section 3.5. This legitimate trace splitting is distinct from the false invariant isomorphism in the Problem.\par
\paragraph*{Koszul criteria.}\mbox{}\par
The star criterion of Theorem 3.6 was checked using the joint kernel of multiplication maps and its alternating part, which is the usual Koszul cycle space. Adding a star diagonal and using the stated \(H^1\)-vanishing supplies a lift; antisymmetrization gives the Koszul boundary in characteristic zero. For Theorem 3.8, the chain ideals give the intersection of adjacent multiplication kernels in \(H^0(L)^{\otimes j}\). Adding the last diagonal and using the stated vanishing proves exactness degree by degree of the corresponding linear free complex. These checks support the criteria in the main connected-projective setting. The full original Green and Inamdar–Mehta papers were not acquired; the historical correction claimed in Remark 3.7 was not independently adjudicated.\par
\paragraph*{Vanishing and endpoint arithmetic.}\mbox{}\par
For \(m\ge n+1\), \(L_m\) is globally generated, hence nef: Kodaira vanishing gives the required Castelnuovo–Mumford regularity relative to the finite map \(\alpha\). Consequently \[L_m^{\,k}\otimes\mathcal O(lA)
 =\omega_X\otimes\bigl(L_m^{\,k-1}\otimes\mathcal O((m+l)A)\bigr)\] has vanishing higher cohomology when \(m+l\ge1\). This supplies the implicit nefness behind Lemma 5.2.\par
The exterior-power evaluation sequence yields both inductions in Lemma 5.3. In Lemma 5.4, if all nonzero cohomology degree is on the first factor, then \(i\ge d-r+2\) and \(i\le n\) imply \(d\le n+r-2<m\). Lemmas 5.5–5.6 force either a vanishing factor or an impossible strictly increasing exponent sequence; deletion of zero-degree endpoint factors preserves the needed inequality. The induction on zero complex degrees in Lemma 5.7 shifts the required cohomology degree by at least \(j-1\), exactly matching its bound. Theorem 5.8 consequently has the thresholds needed by the star and chain criteria. These checks do not replace the subsequent diagonal-lifting obligation.\par
There are minor endpoint clarifications. Lemma 5.1 includes \(s=1\), for which the empty tensor is \(\mathcal O\) in degree zero and is not an exact augmented complex; restricting the assertion to \(s\ge2\) covers the main uses. The index range in the first line of Lemma 5.6’s proof should run through \(r\), not \(n\). These local issues are not promoted to failures of the main vanishing theorem.\par
\paragraph*{Higher lifting and powers.}\mbox{}\par
The chain induction has a valid combinatorial stabilization: when replacing the edge \((s,s+1)\), act on factor \(s\). Earlier edges are enlarged \(\Gamma\)-conditions invariant under the action, while later ordinary diagonals avoid that factor. Thus a supposed failure merely because one factor meets two neighboring chain edges is not established. The ramified section-gluing problem remains.\par
The final proof must cover every \(k\ge1\) required in Theorems 3.6 and 3.8. Sections 5.4–5.5 display \(F\) with first powers, whereas Lemma 5.10 and the criteria involve unequal powers. Tensoring a sheaf map with a line bundle does not automatically preserve surjectivity on global sections. A replacement lifting argument must explicitly handle these powers as well as mixed-union compatibility; this is a specified remaining obligation, not a separate claimed counterexample.\par
\section*{Conclusion and state}
The report identifies one demonstrated core proof defect and gives both local and global evidence under the main hypotheses. The successful enlarged-diagonal vanishing and the positive unrestricted test case are retained as counterevidence. General lifting at ramification remains to be justified.\par
Item | Assessment | Repair verification | Adoption | Remaining work\par
Invariant gluing, p. 15 | Major | No general repair certified | Not requested | Regular compatible lifts\par
Enlarged vanishing | Positive bounded check | Indices and thresholds checked | None | Does not imply lifting\par
Unused endpoints | Minor clarification | \(s\ge2\); index through \(r\) | Not requested | Correct local wording\par
All powers and unions | Needs Evidence | Test case succeeds | None | General replacement proof\par
Initial review | Pending | No independent acceptance | None | Independent acceptance\par
\section*{Scope and references}
All six actual supplied input files were checked by byte count and SHA256. The review records distinguish complete text reading, selected source-image inspection, bounded proof checks, exact reviewer calculations, and unverified external dependencies. Manifest parsing and hashing are not counted as scientific verification. No author code was run.\par
Every final report page was inspected. The external controller owns the supplied Chat/GPT-6/Pro UI observation; backend identity was not independently verified. Scientific adjudication is pending, human reviews are empty, and author adoption was not requested.\par
9 P. Bangere and J. Lacini, Syzygies of adjoint linear series on projective varieties, arXiv:2302.02517v1, 6 February 2023. https://arxiv.org/abs/2302.02517v1.\par
\end{document}
